\section{Research Projects \& Publications}


    \begin{samepage}
        \begin{twocolentry}{
            \href{https://conf.researchr.org/home/icse-2026}{ICSE 2026}
        }
            \textbf{Testing Deep Learning Libraries via Neurosymbolic Constraint Learning}
        \end{twocolentry}

        \details{\vspace{0.05 cm}}
        \vspace{0.05 cm}

        \begin{onecolentry}
            \mbox{\textbf{\textit{M M Abid Naziri}}\textsuperscript{*}, Shinhae Kim\textsuperscript{*}, Feiran Qin, Saikat Dutta, Marcelo d'Amorim}\\
            {\footnotesize \textsuperscript{*}Equal contribution}

            % \vspace{0.10 cm}

            \details{\begin{highlights}
                \item Designed and implemented \textbf{Centaur}, a neurosymbolic constraint learning tool for DL APIs to generate valid test inputs.
                \item Uncovered \textbf{23} new bugs in \href{https://github.com/tensorflow/tensorflow}{Tensorflow} and \href{https://github.com/pytorch/pytorch}{PyTorch} (\textbf{11} confirmed).
            \end{highlights}}

            % \vspace{0.10 cm}

            \textbf{Venue:} \href{https://conf.researchr.org/details/icse-2026/icse-2026-research-track/272/Testing-Deep-Learning-Libraries-via-Neurosymbolic-Constraint-Learning}{IEEE/ACM International Conference on Software Engineering (ICSE 2026)} \href{https://arxiv.org/pdf/2601.15493}{[pdf]}
    \end{onecolentry}
    \end{samepage}

    \details{\vspace{0.05 cm}}
    \vspace{0.15 cm}

    \begin{samepage}
        \begin{twocolentry}{
            \href{https://conf.researchr.org/home/icse-2026}{ICSE 2026}
        }
            \textbf{Misbehavior Forecasting for Focused Autonomous Driving Systems Testing}
        \end{twocolentry}

        \details{\vspace{0.05 cm}}
        \vspace{0.05 cm}

        \begin{onecolentry}
            \mbox{\textbf{\textit{M M Abid Naziri}}, Stefano Carlo Lambertenghi, Andrea Stocco, Marcelo d'Amorim}

            % \vspace{0.10 cm}

            \details{\begin{highlights}
                \item Developed \textbf{Foresee}, a tool to enhance simulation-based testing of self-driving software by identifying and analyzing near-miss scenarios, improving failure detection efficiency and effectiveness.
                \item \textbf{Foresee} exposes \textbf{128\%} more failures than state-of-the-art failure predictors.
                \item \textbf{Foresee} also enhances the capability of an existing technique \href{https://github.com/dk-kling/drivefuzz}{DriveFuzz} by upto \textbf{94\%}.
            \end{highlights}}

            % \vspace{0.10 cm}

            \textbf{Venue:} \href{https://conf.researchr.org/details/icse-2026/icse-2026-research-track/301/Misbehavior-Forecasting-for-Focused-Autonomous-Driving-Systems-Testing}{IEEE/ACM International Conference on Software Engineering (ICSE 2026)} \href{https://arxiv.org/pdf/2512.18823}{[pdf]}
    \end{onecolentry}
    \end{samepage}

    \details{\vspace{0.05 cm}}
    \vspace{0.15 cm}

    \begin{samepage}
        \begin{twocolentry}{
            \href{https://conf.researchr.org/track/issta-2025/issta-2025-tool-demonstrations}{ISSTA 2025}
        }
            \textbf{BugsInDLLs: A Database of Reproducible Bugs in Deep Learning Libraries to Enable Systematic Evaluation of Testing Techniques}
        \end{twocolentry}

        \details{\vspace{0.05 cm}}
        \vspace{0.05 cm}

        \begin{onecolentry}
            \mbox{\textbf{\textit{M M Abid Naziri}}, Aman Kumar Singh, Benjamin Wu, Feiran (Alex) Qin, Saikat Dutta, Marcelo d'Amorim}
            % \mbox{\textbf{\textit{M M Abid Naziri}}}, \mbox{Stefano Carlo Lambertenghi}, \mbox{Andrea Stocco}, \mbox{Marcelo d'Amorim}

            % \vspace{0.10 cm}

            \details{\begin{highlights}
                \item Developed an open-source dataset \textbf{BugsInDLLs} with \textbf{112} reproducible bugs across three popular DL libraries: \href{https://github.com/jax-ml/jax}{JAX}, \href{https://github.com/tensorflow/tensorflow}{Tensorflow}, and \href{https://github.com/pytorch/pytorch}{PyTorch}.
                \item Added support for benchmarking other fuzzing tools with an included integration of \href{https://github.com/ise-uiuc/FreeFuzz}{FreeFuzz}.
            \end{highlights}}

            % \vspace{0.10 cm}

            \textbf{Venue:} \href{https://dl.acm.org/doi/pdf/10.1145/3713081.3731739}{International Symposium on Software Testing and Analysis (ISSTA 2025: Tool Demos)}
            \href{https://abidnazirisami.github.io/papers/bugsindlls.pdf}{[pdf]}
        \end{onecolentry}
    \end{samepage}

    \details{\vspace{0.05 cm}}
    \vspace{0.15 cm}

    \begin{samepage}
        \begin{twocolentry}{
            \href{https://conf.researchr.org/home/ast-2026}{AST 2026}
        }
            \textbf{Improving Deep Learning Library Testing with Machine Learning}
        \end{twocolentry}

        \details{\vspace{0.05 cm}}
        \vspace{0.05 cm}

        \begin{onecolentry}
            \mbox{Facundo Molina, \textbf{\textit{M M Abid Naziri}}, Feiran (Alex) Qin, Alessandra Gorla, Marcelo d'Amorim}

            \details{\begin{highlights}
                \item Applied ML classifiers to infer the validity of an input before executing to increase efficiency
                of DL Library fuzzing tools.
                \item The classifiers achieved an accuracy of \textbf{91}\% in predicting input validity.
                \item They also improved an existing fuzzing tool called \href{https://github.com/shijy16/ACETest}{ACETest} by increasing its validity
                ratio from \textbf{29}\% to \textbf{61}\%.
            \end{highlights}}

            % \vspace{0.10 cm}

            \textbf{Venue:} \href{https://conf.researchr.org/details/ast-2026/ast-2026-papers/4/Improving-Deep-Learning-Library-Testing-with-Machine-Learning}{ACM/IEEE International Conference on Automation of Software Test (AST 2026)} \href{https://abidnazirisami.github.io/papers/specs4freeDLL.pdf}{[pdf]}
    \end{onecolentry}
    \end{samepage}

    \details{\vspace{0.05 cm}}
    \vspace{0.15 cm}


    \begin{samepage}
        \begin{twocolentry}{
            \href{https://conf.researchr.org/home/icst-2026}{ICST 2026}
        }
            \textbf{Evaluating the Effectiveness of Coverage-Guided Fuzzing for Testing Deep Learning Library APIs}
        \end{twocolentry}

        \details{\vspace{0.05 cm}}
        \vspace{0.05 cm}

        \begin{onecolentry}
            \mbox{Feiran Qin, \textbf{\textit{M M Abid Naziri}}, Hengyu Ai, Saikat Dutta, Marcelo d'Amorim}

            % \vspace{0.10 cm}

            \details{\begin{highlights}
                \item Conducted a study to assess the effectiveness of coverage guided fuzzing to test Deep Learning Library APIs.
                \item Reported \textbf{42} new bugs in \href{https://github.com/tensorflow/tensorflow}{Tensorflow} and \href{https://github.com/pytorch/pytorch}{PyTorch} (\textbf{8} Fixed).
            \end{highlights}}

            % \vspace{0.10 cm}

            \textbf{Venue:} \href{https://conf.researchr.org/home/icst-2026}{IEEE International Conference on Software Testing, Verification and Validation (ICST 2026)} \href{https://arxiv.org/pdf/2509.14626}{[pdf]}
    \end{onecolentry}
    \end{samepage}
